\documentclass[aspectratio=169]{beamer}
\usetheme{Madrid}
\usepackage{booktabs}

\definecolor{BrandBlue}{HTML}{0F2744}
\definecolor{BrandTeal}{HTML}{0D9488}
\definecolor{BrandLight}{HTML}{EAF4F4}

\setbeamercolor{title}{fg=white,bg=BrandBlue}
\setbeamercolor{frametitle}{fg=BrandBlue}
\setbeamercolor{structure}{fg=BrandTeal}
\setbeamercolor{block title}{fg=white,bg=BrandBlue}
\setbeamercolor{block body}{bg=BrandLight}
\setbeamertemplate{navigation symbols}{}

\title{Korea Income \& Welfare}
\subtitle{Education, income, and model benchmarking on Korean welfare microdata}
\author{Santiago Torrado}
\date{}

\begin{document}

\begin{frame}
  \titlepage
\end{frame}

\begin{frame}{Why this project matters}
  \begin{itemize}
    \item This case study explores a classic socioeconomic question with a portfolio-ready analytics workflow.
    \item The business and policy angle is simple: how much signal does education carry when explaining income?
    \item The project is useful for analytics, public policy, labor-market research, and applied economics positioning.
  \end{itemize}

  \begin{block}{Dataset}
    92,857 observations, 14 raw variables, individual-level welfare and income data from Korea.
  \end{block}
\end{frame}

\begin{frame}{Analytical question and hypotheses}
  \begin{block}{Core question}
    How much predictive value does education provide when estimating individual income?
  \end{block}

  \begin{itemize}
    \item \textbf{H0:} Education level does not have a meaningful relationship with income.
    \item \textbf{H1:} Education level has a meaningful positive relationship with income.
  \end{itemize}

  \begin{block}{Professional framing}
    This is not presented as a causal claim. The goal is to measure signal, compare model behavior, and communicate findings honestly.
  \end{block}
\end{frame}

\begin{frame}{Workflow}
  \begin{enumerate}
    \item Load and profile the welfare microdata.
    \item Map coded fields into readable categories.
    \item Convert income from thousand KRW to USD and derive age.
    \item Remove sparse columns and filter outliers with the IQR rule.
    \item Benchmark three models using education as the predictive feature:
    \begin{itemize}
      \item Linear Regression
      \item Polynomial Regression
      \item Random Forest Regressor
    \end{itemize}
  \end{enumerate}
\end{frame}

\begin{frame}{Key findings}
  \begin{itemize}
    \item Education shows a positive correlation with income: \textbf{r = 0.404}.
    \item Family size also shows a strong positive correlation: \textbf{r = 0.420}.
    \item Age shows a negative association in the transformed sample: \textbf{r = -0.356}.
    \item The visual analysis supports a steady increase in average income across education tiers.
  \end{itemize}

  \begin{block}{Interpretation}
    Education clearly matters, but it does not explain income on its own. That is a stronger portfolio conclusion than overclaiming model power.
  \end{block}
\end{frame}

\begin{frame}{Model benchmark}
  \small
  \begin{tabular}{lrrr}
    \toprule
    Model & MAE (USD) & RMSE (USD) & $R^2$ \\
    \midrule
    Linear Regression & 1300.23 & 1724.17 & 0.327 \\
    Polynomial Regression & 1299.25 & 1720.02 & 0.330 \\
    Random Forest & 1289.11 & 1713.79 & 0.335 \\
    \bottomrule
  \end{tabular}

  \vspace{0.7cm}

  \begin{block}{Takeaway}
    Random Forest performs best, but only marginally. The modest $R^2$ values show that a single-feature benchmark captures only part of the income process.
  \end{block}
\end{frame}

\begin{frame}{Conclusions and next steps}
  \begin{itemize}
    \item The project supports the alternative hypothesis: education is meaningfully associated with income.
    \item The result is directionally strong, but not sufficient for a high-power predictive model.
    \item The current notebook works best as an interpretable benchmark, not as a production model.
  \end{itemize}

  \begin{block}{Next iteration}
    Add multivariate features, proper categorical encoding, feature importance analysis, and exportable charts for recruiting and LinkedIn use.
  \end{block}
\end{frame}

\end{document}
